\begin{figure}[H]
\centering
\makebox[\textwidth][c]{%
\begin{tikzpicture}[
    scale=0.85,
    transform shape,
    actor/.style={circle, draw=black, fill=white, line width=0.8pt, minimum size=0.7cm},
    usecase/.style={ellipse, draw=black, fill=white, line width=0.55pt, minimum width=3.2cm, minimum height=0.6cm, font=\scriptsize, align=center},
    system/.style={rectangle, draw=black!50, dashed, fill=white, line width=0.4pt, inner sep=8pt, rounded corners=3pt}
]

\node[system, minimum width=12.5cm, minimum height=9cm] (system) at (0,0) {};
\node[font=\bfseries\small] at (0,4.1) {Phân rã: Quản lý hàng hóa \& kho};

\node[usecase] (uc15) at (-3.2,2.5) {UC15: Thêm hàng hóa};
\node[usecase] (uc16) at (-3.2,1.3) {UC16: Sửa hàng hóa};
\node[usecase] (uc17) at (-3.2,0.1) {UC17: Ngừng hàng hóa};
\node[usecase] (uc18) at (-3.2,-1.1) {UC18: Tìm kiếm hàng hóa};

\node[usecase] (uc19) at (3.2,2.5) {UC19: Thêm kho};
\node[usecase] (uc20) at (3.2,1.3) {UC20: Sửa kho};
\node[usecase] (uc21) at (3.2,0.1) {UC21: Tìm kiếm kho};

\node[actor, label=below:{\scriptsize\bfseries Nhân viên kho}] (nvk) at (-7,0.7) {};
\node[actor, label=below:{\scriptsize\bfseries Quản lý kho}] (qlk) at (7,0.7) {};

\foreach \u in {uc15,uc16,uc17,uc18} {
  \draw (nvk) -- (\u);
  \draw (qlk) -- (\u);
}
\draw (qlk) -- (uc19);
\draw (qlk) -- (uc20);
\draw (qlk) -- (uc21);
\draw (nvk) -- (uc21);

\end{tikzpicture}
}
\caption{Sơ đồ use case phân rã -- Hàng hóa và kho}
\label{fig:usecase-item-management}
\end{figure}
